\section{Reconstruction, sharp dyadic displacement, and isometry}\label{sec:dyadic}
\subsection{Odd reconstruction separates identity from precision}
\contract{The odd Euler reconstruction is an exact identity over $\mathbb Z$ for every odd length $p$, even when that length is composite. The prime-only corollary is a specialization.}
\excerpt{Odd-length Euler reconstruction}{odd_euler}
The series identity underlying this theorem multiplies a rescaled squared reciprocal series by the even part of the source series. It then extracts exponential coordinates, proves odd terms vanish, and changes from arbitrary indices to even indices. The coefficient $\binom{2n}{2m}$ and the factor $p^{2m}$ come from that exact coordinate calculation. The source separately proves the signed zigzag identity for the inverse coefficient; the sign is not inferred from a plotted prefix.

\contract{Truncation includes $m=0,\ldots,\min(M,n)$. Every omitted summand has the stated even weight.}
\excerpt{Integral odd-length tail divisibility}{odd_tail}
\walk{The full sum is split into retained and omitted indices. For an omitted $m$, $m\ge M+1$, so its factor $p^{2m}$ contains $p^{2M+2}$. The tail theorem is therefore an integer divisibility statement. If $M\ge n$, there are no omitted terms and the truncation equals the full value exactly. No convergence question arises, because each row has only $n+1$ terms.}

\subsection{The dyadic transform has its own coefficients}
\contract{The dyadic Euler reconstruction uses its specifically constructed dyadic source and integral transform coefficients. It is not obtained by substituting $p=2$ into a theorem whose hypothesis is \code{Odd p}.}
\excerpt{Dyadic reconstruction identity}{dyadic_reconstruct}
The dyadic coefficient analysis includes the valuation of the transform coefficients and of the source terms. Its sharper truncation bound is
\[
2^{4M+3}\mid E_{2n}-T_M(n)
\]
in the p-adic integral coordinates. That exponent comes from the omitted index bound and the valuation of the dyadic coefficient $\kappa_m$. The first omitted term is separately analyzed to show sharpness where the stated row/depth hypotheses apply.

\contract{The displayed uniform sharp tail bound holds for all natural $n,M$. Full-depth truncations have zero error. A nonzero exact valuation for the first omitted error has its own boundary hypotheses.}
\excerpt{Sharper dyadic tail}{sharp_tail}
\walk{The proof splits the finite reconstruction sum by the predicate $m\le M$. In the complement, $m>0$, allowing the coefficient divisibility theorem. The exponent comparison $4M+3\le4m-1$ supplies the uniform bound. The calculation is a universal finite-sum proof; checking numerical rows is independent bounded evidence about examples, not the reason this theorem holds.}

\subsection{Dominance gives an exact valuation}
A lower bound on every difference does not itself determine a least period. The sharp modules isolate a nonzero leading Newton contribution whose valuation is smaller than all remaining contributions. In a nonarchimedean field, that strict separation forces the whole sum to have the leading term's valuation. A tie in the lowest valuations could allow cancellation, so the strict tail estimate is essential.

\contract{For each of the four normalized dyadic addresses $j<4$ and each positive natural displacement $T$, the difference is nonzero and its valuation is $v_2(T)+b_j$, where $b_j$ is the explicitly defined prefix bonus.}
\excerpt{Exact normalized displacement}{sharp_displacement}
\walk{The endpoint supplies four inputs to the generic Newton dominance theorem: first nonvanishing coefficient, its exact valuation, the strict tail-precision bound, and positivity of $T$. The conjunction records nonzero difference as well as valuation equality. At $T=0$ the difference is zero, so the exact nonzero statement would be false; the theorem excludes that case.}

The normalized bonus and the unnormalized bonus are distinguished. The latter incorporates the relevant factorial normalization. The exact-period classification is stated using the actual unnormalized sequence and its own bonus $c_j$.

\contract{A period in the source is a positive displacement. At precision $r$, the four-address unnormalized packet has period $T$ exactly when $T>0$ and $2^{\max(r-c_j,0)}\mid T$. Natural subtraction in the source expresses the maximum with zero.}
\excerpt{Every dyadic period classified}{least_period}
\walk{Exact displacement reduces a universal congruence at every base row to an inequality $v_2(T)+c_j\ge r$. The generic classification translates that inequality into divisibility. If $r\le c_j$, the exponent is zero and every positive displacement is a period. The least positive one is then one. If $r>c_j$, the least period is the indicated power of two. Zero remains excluded by the period predicate.}

\subsection{The normalized packet extends to a p-adic isometry}
\contract{For $j<4$, the appropriately rescaled rational packet extends uniquely to an isometric equivalence $\mathbb Z_2\simeq\mathbb Z_2$ agreeing with its values at natural integers.}
\excerpt{Unique isometry extension}{isometry}
\walk{The source first constructs the isometry and proves agreement on natural inputs. Uniqueness then applies the theorem that continuous extensions agreeing on the dense natural subset agree everywhere. The isometric-equivalence conclusion includes surjectivity, so a merely distance-preserving embedding would not suffice. The rescaling and basepoint translation are part of the packet definition; the raw unnormalized sequence has a different image ball.}

The concrete normalization is
\[
 (b_0,b_1,b_2,b_3)=(1,1,4,2),\qquad
 (c_0,c_1,c_2,c_3)=(1,1,5,3),
\]
where $b_j$ are normalized-coefficient bonuses and $c_j$ are original-coefficient bonuses. Writing $U_{4,j}$ for the unnormalized packet, the isometric packet is
\[
 P_j(n)=\frac{U_{4,j}(n)-j!}{2^{c_j}}.
\]
Its extension maps the unit ball onto the unit ball. Undoing that formula shows the original unnormalized extension has image $j!+2^{c_j}\mathbb Z_2$.

This distinction explains the manuscript's affine-image consequence: undoing the normalization sends the unit ball to a specified affine residue ball. Neither the period theorem nor the isometry endpoint proves an unrestricted arithmetic statement for all moduli. Their exact hypotheses identify the four-address dyadic family being closed.
